% Conclusion
\section{Conclusion}
\label{sec:conclusion}

We presented a continuous-space multi-agent path finding system that combines EPIBTShield, a priority-based velocity conflict resolution algorithm, with a rectified-flow GNN that provides spatially informed preferred velocities. The two components are cleanly separable: EPIBTShield alone (Straight+EPIBTShield) already delivers a $60.6$ percentage-point improvement over ORCA in dense random maps by resolving the symmetric deadlocks that undermine velocity-obstacle methods under high agent density. Adding the rectified-flow GNN provides a further additive gain of $10.0$ percentage points — most pronounced at $N{=}100$ where global coordination reduces the waiting burden on the shield.

\textbf{EPIBTShield} generalizes discrete PIBT to continuous domains via a scored candidate-velocity set with 19 candidates per agent, strict priority ordering, and recursive backtracking up to depth 3. Unlike ORCA, it never requires a non-empty feasible set: agents that cannot find a collision-free forward velocity commit the zero velocity and retry on the next step, avoiding the pathological behavior where ORCA agents oscillate in symmetric deadlocks for hundreds of steps. \textbf{The rectified-flow GNN} is trained in 200k gradient steps on EECBS expert demonstrations using a simple MSE flow-matching loss, and produces preferred velocities in approximately 8\,ms for $N{=}100$ agents with only 3 Euler integration steps at inference.

The system's primary open challenge is OOD generalization of the flow component. On warehouse maps, the learned preferred velocity is frequently perpendicular to the corridor axis, causing EPIBTShield to commit zero velocities and stall agents. Straight+EPIBTShield, which uses only the geometric direction to goal, degrades gracefully under the same distribution shift. This result cleanly isolates the learned component as the generalization bottleneck.

Future work includes: (1) diversifying the EECBS training distribution to include structured, warehouse, and maze-style maps; (2) integrating a receding-horizon multi-step lookahead into EPIBTShield to reduce collisions from agents that arrive at already-committed positions; (3) replacing the discrete candidate set with a continuous feasibility projection under the priority-ordered commitment constraints; and (4) validating the system on physical multi-robot hardware.
